\section{Source domains, spectral tests, and model comparisons}\label{app:comparisons}
\subsection{Completion is not a physical source prescription}
On the same compact boundary manifold take two smooth
positive gauge-invariant weights $w_0,w_1$. Then
\begin{equation}
U_{01}F=\sqrt{w_0/w_1}\,F
\end{equation}
is unitary between their covariant \Ltwo{} spaces and
intertwines every bounded boundary multiplication.
Smooth positive multiplication and elliptic Sobolev
equivalence also make it a bounded isomorphism of every
fixed positive-order source space $\cK_r$.
Both assertions follow directly by cancellation in
the norm and bounded smooth multiplier estimates.
Thus bare Hilbert occurrence, covariance, and local
regularity bounds do not distinguish these Wilson
weights. The distinguished vector is changed:
$U_{01}I=\sqrt{w_0/w_1}I$. Relations with that vector
or a fixed electric operator are generally not preserved
unless conjugated. Nor is an exact finite-word loop
algebra preserved by arbitrary smooth weight multiplication.
This is why an independently specified state-dependent
source rule is essential.

For clarity, the countable core in
Theorem~\ref{thm:compact-occurrence} can be made explicit:
choose a dense sequence in each rational-endpoint
compact support interval, include the three anchors,
and close over $\Q(\ii)$ under rational translations,
reflection, and rational modulations $e^{\ii qx}$.
Apply $\T$ to that profile space. Partitions of unity
for separated supports are applied to the profiles
h,k, not to the neutral tests $\T h,\T k$, since
arbitrary multiplication need not preserve neutrality.
Continuity then extends rational translations and
modulations to real ones.

For an already continuous source J, pointwise
$\norm{J[\lambda^{-2}\T(e^{\ii\lambda x}h)]}^2
=O_h(\log(2+|\lambda|))$ for \emph{every} h on a fixed
support interval implies a common seminorm bound.
Divide these continuous linear maps by
$\sqrt{\log(2+|\lambda|)}$ and apply the uniform
boundedness principle on the Fr\'echet test space.
Continuity in $\lambda$ controls bounded parameter
intervals. This reasoning does not apply when the
pointwise premise is known only on a dense list.

Membership in $\cK_r$ alone gives Sobolev regularity,
not a bounded or smooth observable. If $d=\dim M$,
$r>d/2$ suffices for continuous bounded representatives;
smoothness needs bounds of all orders or a separate
smoothing statement. \Ltwo{} completion suffices for mixed
norms and bounded multiplications, but arbitrary products,
finite loop polynomials, and unbounded insertions have
separate domain obligations.

\subsection{The generator induced by an exact pairing}
If an exact source J exists, stationarity of Q makes
$Jf\mapsto JU_tf$ an isometry on its range. It is
well defined on null vectors and extends to a strongly
continuous unitary group on the source closure.
No independently named YM generator is required for
this conclusion.

Positivity then implies RH, and \eqref{eq:zero-pairing}
identifies the completed source space with
$\ell^2(\{\gamma\},m_\gamma)$ through Fourier evaluation.
Its image is dense: a sequence orthogonal to all
$\widehat f(\gamma)$, $f=\T h$, defines a tempered
point-supported spectral distribution by Cauchy--Schwarz
and zero counting. Its inverse Fourier distribution
annihilates $\T\cD$; multiplication by
$\tau^2+1/4$ then forces the original distribution
to vanish. Translation has one eigenline per
distinct ordinate, with $m_\gamma$ a norm weight.
Multiplicity of a zero does not require a Jordan
block or $m_\gamma$ independent eigenvectors.

This generator is unbounded in both spectral
directions and cannot be periodic. One can check
these requirements even before invoking the exact
zero model: the modulations of
Proposition~\ref{prop:highfreq} obey
$Q[f_\lambda']/Q[f_\lambda]\sim\lambda^2$,
and reflection sends the generator to its negative.
If $U_P$ acted identically for some $P>0$, choose h
in a support shorter than P. The neutral difference
$f_\lambda-U_Pf_\lambda$ would have zero source norm,
whereas its disjoint leading profiles give
$Q[f_\lambda-U_Pf_\lambda]
=2\norm h^2\log|\lambda|+O(1)$, a contradiction.
Finite-dimensional generators and fixed periodic
color rotations therefore cannot realize it.

Nor can arithmetic translation be identified with
$e^{-tH_{\rm phys}}$ for a nonnegative Euclidean
Hamiltonian on the source vectors. Exact norm
stationarity for all positive t forces every source
into $\ker H_{\rm phys}$ by the spectral theorem.
Then $JU_tf=Jf$, and differentiation gives $Jf'=0$,
contradicting the high-frequency test. This excludes
that ansatz, not general use of a physical Hamiltonian.

A well-defined generalized spectral source associated
with a self-adjoint G and vector $\eta$ is
\begin{equation}
Jg=\{\widehat g(\lambda)(1+\lambda^2)^{s/2}\}(G)\eta,
\quad
\ip{Jf}{Jg}=
\int\overline{\widehat f(\lambda)}\widehat g(\lambda)
                    (1+\lambda^2)^s\dd\mu_\eta(\lambda).
\end{equation}
For smooth compact tests the combined multiplier is
bounded, with fixed-support test seminorm bounds.
It must not be factored into an unbounded operation
on $\eta$ unless that separate domain is verified.
Multiplication by this spectral weight cannot turn
an absolutely continuous spectral measure into
the required atoms.
For a complete weighted boundary flow, multiplication
by $\sqrt w$ conjugates the half-divergence generator
to the corresponding Haar generator; a smooth density
alone does not create an arithmetic spectrum.

\subsection{Modular, boost, and thermal controls}
The boundary matrix multiplication algebra has the
faithful finite trace $\tau(B)=\int\tr(B)\dd\nu/N_c$.
In its trace GNS space, $S(B)=B^*$ is antiunitary,
so the modular operator is $\Delta=I$. A nonconstant
central density does not supply a nontrivial modular
Hamiltonian in this algebra.

For a proposed continuum wedge alternative, assume
the positive-energy Poincar\'e representation, a
unique vacuum, and the applicable wedge modular
identification with boosts \cite{bw}. A boost
eigenvector has a finite joint momentum spectral
measure invariant under that boost: the phase of
the eigenvector cancels. On $p_\pm=p_0\pm p_1>0$,
the logarithms of these coordinates translate under
boosts. A finite translation-invariant measure on
$\R$ must vanish, by disjoint translates of an
interval. Thus the momentum measure is supported
where $p_+=p_-=0$. The forward-cone condition forces
the full momentum to be zero, and uniqueness of
the vacuum leaves only the vacuum vector.
Hence nonzero arithmetic eigenvalues cannot occur
as ordinary wedge-boost eigenvectors under these
hypotheses. No mass gap is needed for this argument.
The wedge identification for a YM net is itself
conditional; it does not follow from the finite slab.

A thermal Liouville comparison changes the state.
For a discrete Hamiltonian $H|n\rangle=E_n|n\rangle$
with finite Gibbs partition function, put
$r_n=Z^{-1}e^{-\beta E_n}$ and
$L=H_{\rm left}-H_{\rm right}$ on Hilbert--Schmidt
operators. For bounded self-adjoint A and
$v=A\rho_\beta^{1/2}$, direct matrix expansion gives
\begin{equation}
\dd\mu_v(\lambda)=
\sum_{m,n}r_n|A_{mn}|^2\delta_{E_m-E_n}(\dd\lambda),
\qquad \mu_v(-\dd\lambda)=e^{-\beta\lambda}\mu_v(\dd\lambda).
\end{equation}
This detailed balance law is incompatible with an
even measure away from zero. The centered choice
$v_c=\rho_\beta^{1/4}A\rho_\beta^{1/4}$ gives instead
weights $\sqrt{r_mr_n}|A_{mn}|^2$, which are even.
Schatten H\"older gives $\norm{v_c}_{\rm HS}\le\norm A$.
Consequently ordinary orbit smearing is still locally
bounded and fails Proposition~\ref{prop:highfreq}.
An independently justified generalized source would
need its actual spectrum and domains analyzed.
In the vacuum, $[H,A]\Omega_{\rm vac}=HA\Omega_{\rm vac}$;
this is not a two-sided thermal Liouville spectrum.

Arithmetic quotient representations such as Meyer's
\cite{meyer} supply a different comparison. Their
virtual character is poles minus zeros, so its sign
on pole-neutral tests is opposite the desired zero
contribution. Derivative-evaluation jets occur in
the continuous dual and are not automatically Hilbert
vectors. If a proposal insists on realizing all such
jets faithfully with a positive norm and unitary
translation, a nontrivial nilpotent Jordan block is
impossible: its exponential has polynomially growing
orbits. That ansatz would also require simple zeros.
The Weil norm does not require simplicity, since
it uses multiplicities as positive weights.

\subsection{What topology and the two-dimensional comparison do and do not change}
The winding/electric local bound used a finite compact
link manifold, smooth positive state, and a simple
contour with distinct links, not its contractibility.
Changing to another such finite graph or a simple
noncontractible cycle preserves the same type of
obstruction. More generally, a limit of pairings
$B_\alpha[f]\le C_I\norm f_2^2$ with a common
fixed-support constant obeys the same bound.
Finite sums, and Hilbert direct sums with summable
squared bounds, do not evade it. A continuum or
singular-state limit could help only if that bound
loses uniformity, while existence and positivity
of the new source still need proof.

The bare character Mellin law also fails this test:
\[
J_\chi f=\sum_{n\ge1}n^{-1/2}\widehat f(\log n)\chi_n,
\qquad
\ip{J_\chi f}{J_\chi g}
=\sum_{n,m}\frac{\overline{\widehat f(\log n)}
\widehat g(\log m)}{\sqrt{nm}}G_{nm}.
\]
The coefficient form is bounded by
$C\int(1+x^2)|f(x)|^2\dd x$ by logarithmic sampling
and \eqref{eq:gram}. Character dressing improves
the winding operator algebra, but not this bare law.

For comparison only, two-dimensional SU(2) YM has
Haar character basis and, in the chosen Casimir units,
cylinder and handle operators
\begin{equation}
P_t\chi_n=e^{-t(n^2-1)}\chi_n,\qquad
H_{\rm hand}\chi_n=n^{-2}\chi_n,\qquad
\mu(\chi_n\otimes\chi_m)=\delta_{nm}n^{-1}\chi_n .
\end{equation}
These are the usual character gluing formulas
\cite[Proposition 5.7]{runkel-szegedy}. On Haar packets
they give exactly
\begin{equation}
N^{2j}H_{\rm hand}^jS_Rf=S_R(e^{-2jx}f),\qquad
P_{t/N^2}S_Rf=e^{t/N^2}S_R(e^{-te^{2x}}f).
\end{equation}
Thus the rescaled actions tested here are local
multipliers, while fixed positive area or a fixed
positive number of handles kills the packets.
The logarithm $L_{\rm hand}=-\tfrac12\log H_{\rm hand}$
satisfies $L_{\rm hand}V_a=V_a(L_{\rm hand}+\log a)$
and $(L_{\rm hand}-R)S_Rf=S_R(xf)$; it supplies the
integer scale, not arithmetic differentiation.
For genus $g\ge2$, the zero-area partition function
is $\sum_n n^{2-2g}=\zeta(2g-2)$. This special-value
identity is not a mixed Weil norm. SO(3) keeps odd
dimensions and changes the dimension Dirichlet
series to $(1-2^{-s})\zeta(s)$. These are changes
of dimension or gauge group, not results in the
fixed four-dimensional state.

Finally the modular-surface scattering coefficient is
$C_{\rm sc}(s)=\Lambda_\xi(2s-1)/\Lambda_\xi(2s)$,
where $\Lambda_\xi(s)=\pi^{-s/2}\Gamma(s/2)\zeta(s)$
\cite[Section 7.3.2]{levitin-strohmaier}.
The entire completion $\xi(s)=s(s-1)\Lambda_\xi(s)/2$
gives the algebraic relation
\begin{equation}
\frac{\xi(z)}{\xi(z+1)}
=\frac{z-1}{z+1}C_{\rm sc}((z+1)/2).
\end{equation}
This is one meromorphic transfer-function comparison
at shift $1/2$, not a YM preparation/evolution/readout
law. The separate project target
$K_\omega(z)=\xi(1/2+z-\omega)/\xi(1/2+z+\omega)$
and its ordinary causal energy balance require
additional physical identities beyond a reflected
Weil pairing. Arithmetic thermal partition functions,
adelic trace formulas, or proposed arithmetic
cohomological pairings likewise do not by themselves
supply those identities.
